\documentclass[12pt,a4paper]{ctexart}

% === 核心包 ===
\usepackage[backend=bibtex,style=numeric,sorting=none]{biblatex}
\usepackage{geometry}
\usepackage{graphicx}
\usepackage{booktabs}
\usepackage{longtable}
\usepackage{array}
\usepackage{calc}
\usepackage{hyperref}
\usepackage{xcolor}
\usepackage{enumitem}
\usepackage{etoolbox}
% longtable 用小号字，缓解 pandoc 表格列宽溢出
\AtBeginEnvironment{longtable}{\small}
% pandoc 输出的列表紧排命令（由 pandoc 生成但未定义）
\providecommand{\tightlist}{%
  \setlength{\itemsep}{0pt}\setlength{\parskip}{0pt}}
% 特殊 Unicode 符号映射（默认拉丁字体缺字）
\usepackage{newunicodechar}
\usepackage{amssymb}
% 星号与运算符
\newunicodechar{⭐}{\ensuremath{\bigstar}}
\newunicodechar{≤}{\ensuremath{\leq}}
\newunicodechar{≥}{\ensuremath{\geq}}
\newunicodechar{≈}{\ensuremath{\approx}}
% 希腊字母（热电主题正文常用）
\newunicodechar{β}{\ensuremath{\beta}}
\newunicodechar{κ}{\ensuremath{\kappa}}
\newunicodechar{σ}{\ensuremath{\sigma}}
% 下标
\newunicodechar{₀}{\textsubscript{0}}\newunicodechar{₁}{\textsubscript{1}}
\newunicodechar{₂}{\textsubscript{2}}\newunicodechar{₃}{\textsubscript{3}}
\newunicodechar{₄}{\textsubscript{4}}\newunicodechar{₈}{\textsubscript{8}}
\newunicodechar{₆}{\textsubscript{6}}\newunicodechar{₇}{\textsubscript{7}}
\newunicodechar{ₑ}{\textsubscript{e}}\newunicodechar{ₗ}{\textsubscript{l}}
\newunicodechar{ₓ}{\textsubscript{x}}
% 上标
\newunicodechar{⁺}{\ensuremath{^{+}}}
\newunicodechar{⁴}{\textsuperscript{4}}
\newunicodechar{₋}{\ensuremath{^{-}}}
\newunicodechar{ᵧ}{\textsuperscript{$\gamma$}}
% 带圈数字 ①–⑧
\newunicodechar{①}{\textcircled{\raisebox{-0.5pt}{1}}}
\newunicodechar{②}{\textcircled{\raisebox{-0.5pt}{2}}}
\newunicodechar{③}{\textcircled{\raisebox{-0.5pt}{3}}}
\newunicodechar{④}{\textcircled{\raisebox{-0.5pt}{4}}}
\newunicodechar{⑤}{\textcircled{\raisebox{-0.5pt}{5}}}
\newunicodechar{⑥}{\textcircled{\raisebox{-0.5pt}{6}}}
\newunicodechar{⑦}{\textcircled{\raisebox{-0.5pt}{7}}}
\newunicodechar{⑧}{\textcircled{\raisebox{-0.5pt}{8}}}
% 罗马数字
\newunicodechar{Ⅱ}{II}
% 状态图标
\newunicodechar{✅}{\checkmark}
\newunicodechar{⏳}{[TBD]}
\newunicodechar{⚠}{\textbf{/!\\}}
\geometry{margin=2.5cm}
\addbibresource{references.bib}

% === 元数据 ===
\title{ 文献调研报告：热电材料 ZT 优化——从材料体系到构效关系的系统综述 }
\date{ 2026-08 }
\author{Pi-Agent（自主文献调研 Agent）}

\begin{document}
\maketitle
\sloppy  % 允许宽松换行，缓解长 URL/DOI 溢出

\section{文献调研报告：热电材料 ZT 优化------从材料体系到构效关系的系统综述}\label{ux6587ux732eux8c03ux7814ux62a5ux544aux70edux7535ux6750ux6599-zt-ux4f18ux5316ux4eceux6750ux6599ux4f53ux7cfbux5230ux6784ux6548ux5173ux7cfbux7684ux7cfbux7edfux7efcux8ff0}

\begin{quote}
\textbf{2026-08-11 量化验证补充}：本报告为阶段一（文献调研）产物。阶段二（构效关系发现） 的量化验证结果见 \texttt{discovery/discovery\_report.md}（量化验证结果汇总章节）与 \texttt{discovery/model\_comparison\_0/1/2/4.md}、\texttt{discovery/symbolic\_0/1/2/4.md}、 \texttt{discovery/cross\_theme\_connections.md}。核心结论：跨材料温度-ZT 单变量建模 无显著标度（负结果，强化 Gap 1），ZT 预测需多描述符（带隙/有效质量/声速）。
\end{quote}

\textbf{主题}: Thermoelectric Materials ZT Optimization \textbf{文献规模}: 209 篇（arxiv 160 + semantic\_scholar 49） \textbf{日期}: 2026-08-02

\begin{center}\rule{0.5\linewidth}{0.5pt}\end{center}

\subsection{1. 执行摘要}\label{ux6267ux884cux6458ux8981}

热电材料实现热-电直接转换，ZT = S²σT/(κₑ+κₗ) 是核心性能指标。本调研覆盖 209 篇文献，系统梳理了室温（Bi₂Te₃）、中温（PbTe、方钴矿）、高温（SnSe、half-Heusler、GeTe）及新兴二维体系的热电材料研究进展。\textbf{核心发现}：

\begin{enumerate}
\def\labelenumi{\arabic{enumi}.}
\tightlist
\item
  \textbf{实验 ZT 纪录}：p 型 SnSe 单晶 ZT≈2.6（923 K）为实验最高；Yb/In 填充方钴矿 ZT≈1.72（773 K）；共振掺杂 half-Heusler ZT≈1.5（980 K）
\item
  \textbf{四条优化主线}：能带工程（收敛/共振掺杂）、声子工程（rattling/纳米结构/软晶格）、双极抑制（带隙工程）、数据驱动（ML/ab initio 输运）
\item
  \textbf{六大研究空白}：ML 预测-实验断层（置信度 0.92）、n 型 SnSe 落后（0.88）、方钴矿复现性（0.75）、双极掺杂温度演化（0.80）、散射相图普适性（0.72）、共振掺杂定量标度（0.70）
\end{enumerate}

\subsection{2. 文献综述}\label{ux6587ux732eux7efcux8ff0}

\subsubsection{2.1 Bi₂Te₃ 基（室温经典）}\label{biux2082teux2083-ux57faux5ba4ux6e29ux7ecfux5178}

Bi₂Te₃ 是室温热电基准材料，面临两大本征挑战：强各向异性与高温双极激发（\cite{TE065}）。优化策略包括点缺陷工程、织构取向、能带隙增大、溶液法/球磨/甩带纳米结构降热导（\cite{TE065}）。低维化（原子五重层薄膜、拓扑绝缘体薄膜）理论预测 ZT 可达 2.0--7.2（\cite{TE062}/\cite{TE063}），但均未实验验证。n 型 Bi₂Te₃ 通过 Se 掺杂实现高性能（\cite{TE070}）。

\subsubsection{2.2 PbTe 基（中温）}\label{pbte-ux57faux4e2dux6e29}

PbTe 是研究最深入的中温体系。\textbf{关键机制}：PbTe 带隙随温度显著增大，且实际带隙在 300--900 K 贴合最优带隙，是 n 型 PbTe 高 ZT 的本质原因（\cite{TE085}）；温度诱导能带收敛+谷间散射主导 p 型输运（\cite{TE090}）；Sr/Na 共掺实验验证能带收敛（\cite{TE105}）。合金化（Ⅱ族）有效降热导（\cite{TE098}）；PbTe 的巨非谐声子散射（\cite{TE097}）与铁电软模附近宽带声子散射（\cite{TE109}）是其本征低热导根源。

\subsubsection{2.3 SnSe 基（高温纪录）}\label{snse-ux57faux9ad8ux6e29ux7eaaux5f55}

SnSe 单晶自 2014 年报道 ZT=2.6 以来成为研究热点（\cite{TE110} 综述）。本征超低热导（0.08 W/mK，\cite{TE126}）源于晶格软化+纳米级 Sn 偏心位移（\cite{TE124}/\cite{TE133}）+强非谐性。\textbf{核心瓶颈}：n 型 SnSe 远落后 p 型（\cite{TE122}），受限于本征 Sn 空位缺陷化学（天然 p 型）与掺杂剂选择；电弧熔炼 n 型多晶已达 ZT=1.8（\cite{TE115}），但 n 型单晶 Pb 掺杂仅 PF=1.2（\cite{TE119}）。应变工程与织构化是新兴调控手段（\cite{TE120}）。

\subsubsection{2.4 Half-Heusler（高温稳定）}\label{half-heuslerux9ad8ux6e29ux7a33ux5b9a}

Half-Heusler 以高热稳定性、无毒著称，主流 ZT≈1（\cite{TE138}/\cite{TE143}）。\textbf{突破}：Hf₀.₃Zr₀.₇CoSn₀.₃Sb₀.₇ 经 \textless1at\% Al 共振掺杂，PF ↑65\%、κ ↓13\%，ZT 0.8→1.5 @980K（\cite{TE146}）；ZrNiSn 基建立了载流子散射相图------化学掺杂同时屏蔽电离杂质/晶界/极性光学声子散射，但对 κₗ 影响可忽略（\cite{TE136}）；FeNb₀.₈Ti₀.₂Sb 的高性能源于本征电子结构（\cite{TE200}）。

\subsubsection{2.5 填充方钴矿（中温工程典范）}\label{ux586bux5145ux65b9ux94b4ux77ffux4e2dux6e29ux5de5ux7a0bux5178ux8303}

方钴矿通过填充原子降低晶格热导。\textbf{2026 年突破}：Yb 填充+In 填充+Te 掺杂实现双带优化+缺陷工程协同，ZT≈1.72 @773K（PF=5.63 mW/mK²，κ=2.23 W/mK），7 对单级器件转换效率 8.2\%@ΔT=380K（\cite{TE040}）。Yb 填充+CoSi₂ 纳米颗粒达 ZT=1.43（\cite{TE039}）。\textbf{机制修正}：X 射线动力学衍射表明 rattling 实为填充物-笼系统整体振动动力学，而非独立低频振动（\cite{TE188}）。\textbf{挑战}：p 型复现性差（ZT 0.67--0.89 波动，\cite{TE048}），Yb 溶解度温度依赖性复杂（\cite{TE035}）。

\subsubsection{2.6 GeTe / 新兴二维体系}\label{gete-ux65b0ux5174ux4e8cux7ef4ux4f53ux7cfb}

GeTe 通过功函数+声阻抗失配界面协同散射达 ZT≈1.93（\cite{TE165}）；极性畸变导致负热膨胀与软声子（\cite{TE163}）。二维体系（Graphdiyne ZT=3.0、Janus ISbTe 应变 ZT=1.31、β-PdXY ZT=0.68、PbSe/PbTe 单层 ZT=5.3）均为理论预测，代表未验证空间。

\subsubsection{2.7 数据驱动与计算}\label{ux6570ux636eux9a71ux52a8ux4e0eux8ba1ux7b97}

\begin{itemize}
\tightlist
\item
  \textbf{ML 预测}：已有 R²=0.90--0.98 的 ZT/热导预测模型（\cite{TE011}/\cite{TE057}），但预测-实验断层显著（\cite{TE057}）
\item
  \textbf{ab initio 输运}：从电子-声子散射率→BTE→Monte Carlo 纳米结构模拟的完整流程（\cite{TE026}）
\item
  \textbf{品质因子方法}：区分本征材料属性与载流子浓度优化的有效质量模型（\cite{TE025}）
\item
  \textbf{LLM 数据库}：大语言模型驱动的热电材料数据库（\cite{TE017}）
\end{itemize}

\subsection{3. 关键材料与性质对比}\label{ux5173ux952eux6750ux6599ux4e0eux6027ux8d28ux5bf9ux6bd4}

{\def\LTcaptype{none} % do not increment counter
\begin{longtable}[]{@{}
  >{\raggedright\arraybackslash}p{(\linewidth - 14\tabcolsep) * \real{0.0896}}
  >{\raggedright\arraybackslash}p{(\linewidth - 14\tabcolsep) * \real{0.0896}}
  >{\raggedright\arraybackslash}p{(\linewidth - 14\tabcolsep) * \real{0.1642}}
  >{\raggedright\arraybackslash}p{(\linewidth - 14\tabcolsep) * \real{0.0896}}
  >{\raggedright\arraybackslash}p{(\linewidth - 14\tabcolsep) * \real{0.1940}}
  >{\raggedright\arraybackslash}p{(\linewidth - 14\tabcolsep) * \real{0.1493}}
  >{\raggedright\arraybackslash}p{(\linewidth - 14\tabcolsep) * \real{0.1343}}
  >{\raggedright\arraybackslash}p{(\linewidth - 14\tabcolsep) * \real{0.0896}}@{}}
\toprule\noalign{}
\begin{minipage}[b]{\linewidth}\raggedright
材料
\end{minipage} & \begin{minipage}[b]{\linewidth}\raggedright
类型
\end{minipage} & \begin{minipage}[b]{\linewidth}\raggedright
ZT（实验）
\end{minipage} & \begin{minipage}[b]{\linewidth}\raggedright
温度
\end{minipage} & \begin{minipage}[b]{\linewidth}\raggedright
PF (mW/mK²)
\end{minipage} & \begin{minipage}[b]{\linewidth}\raggedright
κ (W/mK)
\end{minipage} & \begin{minipage}[b]{\linewidth}\raggedright
主要机制
\end{minipage} & \begin{minipage}[b]{\linewidth}\raggedright
证据
\end{minipage} \\
\midrule\noalign{}
\endhead
\bottomrule\noalign{}
\endlastfoot
SnSe 单晶 (p) & 实验 & 2.6 & 923 K & --- & 0.08 & 晶格软化+超低热导 & \cite{TE113}/126 \\
n 型 SnSe 多晶 & 实验 & 1.8 & 773 K & --- & --- & 电弧熔炼+织构 & \cite{TE115} \\
Yb/In 方钴矿 & 实验 & 1.72 & 773 K & 5.63 & 2.23 & 双带优化+缺陷工程 & \cite{TE040} \\
HfZrCoSnSb+Al & 实验 & 1.5 & 980 K & ↑65\% & ↓13\% & 共振掺杂 & \cite{TE146} \\
Yb 方钴矿+CoSi₂ & 实验 & 1.43 & \textasciitilde800 K & --- & --- & 纳米颗粒 & \cite{TE039} \\
(Hf,Zr)NiSn & 实验 & 1.2 & \textasciitilde900 K & --- & --- & 合金化 & \cite{TE007} \\
GeTe 基 & 实验 & 1.93 & --- & --- & --- & 界面声阻抗失配 & \cite{TE165} \\
Bi₂Te₃ (经典) & 实验 & \textasciitilde1.0 & 300-500 K & --- & --- & 点缺陷+织构 & \cite{TE065} \\
\end{longtable}
}

\textbf{理论预测 ZT（未验证）}：Bi₂Te₃ 薄膜 7.2（\cite{TE062}）、PbSe/PbTe 单层 5.3（\cite{TE096}）、Si-Ge-P 3.6（\cite{TE002}）、Graphdiyne 3.0（\cite{TE009}）、(PbTe)₂ 层 2.9（\cite{TE100}）。

\subsection{4. 研究空白与未来方向}\label{ux7814ux7a76ux7a7aux767dux4e0eux672aux6765ux65b9ux5411}

详见 \texttt{gap\_report.md}（Gap 1--6）。核心优先方向： 1. \textbf{Gap 1（高）}：弥合 ML 预测-实验断层------主动学习+相稳定性筛选闭环（\cite{TE057}） 2. \textbf{Gap 2（高）}：n 型 SnSe 缺陷化学-掺杂-织构协同优化（\cite{TE122}） 3. \textbf{Gap 6（中）}：共振掺杂 DOS 异常-塞贝克增强定量标度律（\cite{TE146}） 4. \textbf{Gap 4（中）}：双极最优掺杂随温度演化的实验验证（\cite{TE082}） 5. \textbf{Gap 3（中）}：方钴矿批次复现性标准协议（\cite{TE048}） 6. \textbf{Gap 5（中）}：载流子散射相图跨体系普适性验证（\cite{TE136}）

\subsection{5. 参考文献（代表性，含 DOI）}\label{ux53c2ux8003ux6587ux732eux4ee3ux8868ux6027ux542b-doi}

{\def\LTcaptype{none} % do not increment counter
\begin{longtable}[]{@{}
  >{\raggedright\arraybackslash}p{(\linewidth - 4\tabcolsep) * \real{0.2667}}
  >{\raggedright\arraybackslash}p{(\linewidth - 4\tabcolsep) * \real{0.4000}}
  >{\raggedright\arraybackslash}p{(\linewidth - 4\tabcolsep) * \real{0.3333}}@{}}
\toprule\noalign{}
\begin{minipage}[b]{\linewidth}\raggedright
ID
\end{minipage} & \begin{minipage}[b]{\linewidth}\raggedright
论文
\end{minipage} & \begin{minipage}[b]{\linewidth}\raggedright
DOI
\end{minipage} \\
\midrule\noalign{}
\endhead
\bottomrule\noalign{}
\endlastfoot
\cite{TE065} & Hong et al., ``Fundamental and Progress of Bi2Te3-based Thermoelectric Materials'', 2018 & 10.1088/1674-1056/27/4/048403 \\
\cite{TE085} & Cao et al., ``Thermally induced band gap increase and high thermoelectric figure of merit of n-type PbTe'', 2019 & 10.1016/j.mtphys.2019.100172 \\
\cite{TE110} & Chen et al., ``High-performance SnSe thermoelectric materials: Progress and future challenge'', 2018 & 10.1016/J.PMATSCI.2018.04.005 \\
\cite{TE136} & Ren et al., ``Establishing the carrier scattering phase diagram for ZrNiSn-based half-Heusler thermoelectric materials'', 2020 & 10.1038/s41467-020-16913-2 \\
\cite{TE040} & Rao et al., ``Dual Band Optimization and Defect Engineering Enable Ultrahigh Thermoelectric Performance in Filled Skutterudites'', 2026 & 10.1002/aenm.202506357 \\
\cite{TE146} & Mitra et al., ``Conventional Half-Heusler Alloys Advance State-of-the-Art Thermoelectric Properties'', 2022 & --- \\
\cite{TE057} & Athar \& Jund, ``Beyond Predicted ZT: Machine Learning Strategies for the Experimental Discovery of Thermoelectric Materials'', 2026 & 10.1016/j.aichem.2026.100113 \\
\cite{TE122} & Abbas et al., ``Challenges and strategies in developing high-performance n-type polycrystalline SnSe thermoelectric materials'', 2025 & 10.1016/j.isci.2025.114025 \\
\cite{TE188} & Valério et al., ``Phonon Scattering Mechanism in Thermoelectric Materials Revised via Resonant X-ray Dynamical Diffraction'', 2020 & 10.1557/mrc.2020.37 \\
\cite{TE025} & Kang \& Snyder, ``Transport property analysis method for thermoelectric materials: material quality factor and the effective mass model'', 2017 & --- \\
\end{longtable}
}

\emph{完整 209 篇文献清单见 \texttt{paper\_summaries.md} 与 \texttt{papers\_te.json}。}

\subsection{证据链}\label{ux8bc1ux636eux94fe}

\begin{quote}
本章节由 \texttt{scripts/inject\_evidence\_chain.py} 依据 discovery 产物自动生成，保证赛题红线 1「每个结论都能指回具体文献或数据库记录」在基本任务报告层闭环。 生成时间：2026-08-03 18:14｜假设总数：5｜证据条目总数：27 数据源：discovery/hypotheses.json（必需）、discovery/discovery\_report.json（可选） 回查路径：survey\_report.md → discovery/hypotheses.json（evidence\_chain）→ discovery/discovery\_report.md（Evidence Chain）→ 检索缓存 search\_results.json / 人工核实。 状态说明：「✅ 已溯源」= 编号证据在 discovery 产物中存在且未被引用审计标记；「⚠ 需人工核对」= reference\_audit 标记该编号不可追溯，须人工确认。
\end{quote}

\begin{center}\rule{0.5\linewidth}{0.5pt}\end{center}

\subsubsection{假设 1｜纳米晶尺寸调控 Si-Ge-P 超饱和固溶体的晶格热导率与 ZT（hypo\_1）}\label{ux5047ux8bbe-1ux7eb3ux7c73ux6676ux5c3aux5bf8ux8c03ux63a7-si-ge-p-ux8d85ux9971ux548cux56faux6eb6ux4f53ux7684ux6676ux683cux70edux5bfcux7387ux4e0e-zthypo_1}

\textbf{结论简述}：针对 ML 预测高 ZT 材料缺乏实验验证的断层，提出以超饱和 Si-Ge-P 固溶体为模型体系，系统改变球磨时间与低温烧结条件，获得晶粒尺寸在 5--50 nm 的系列样品。假设纳米晶界主要散射中长波声子，而 Fe/P 共掺杂调整电子结构，使功率因子得以保留；因此存在最优晶粒尺寸使晶格热导率大幅下降而载流子迁移率未严重恶化。该假设可通过变温 Hall 效应、电导率/Seebeck 系数和激光闪射法热导率联合测试直接验证，从而在实验中逼近 ML 预测的 ZT 值。

\textbf{证据编号列表}： - \texttt{\cite{TE002}}（论文编号） - \texttt{\cite{TE057}}（论文编号） - \texttt{\cite{TE009}}（论文编号） - \texttt{{[}Novelty\ Verification{]}\ Overlap:\ none\ \textbar{}\ Novelty:\ 0.800\ (was\ 0.800)\ \textbar{}\ Queries:\ 3\ \textbar{}\ Results:\ 4}（新颖性查重）

\textbf{来源归属}： - 论文证据：\texttt{\cite{TE002}}、\texttt{\cite{TE057}}、\texttt{\cite{TE009}}（源自 hypotheses.json → evidence\_chain） - 新颖性查重：新颖性 0.800；查询 3 次；结果 4 条 - 数据库记录：无（hypotheses.json 中 external\_validation 为空）

\textbf{可追溯状态}：✅ 已溯源 ------ 证据编号在 discovery 产物（hypotheses.json → evidence\_chain、discovery\_report.md → Evidence Chain）中可逐层回查，且未被引用审计标记；论文编号对应的真实文献最终确认请回查检索缓存或人工复核。

\begin{center}\rule{0.5\linewidth}{0.5pt}\end{center}

\subsubsection{假设 2｜卤素掺杂位点与 n 型 SnSe 载流子浓度及 ZT 的关联（hypo\_2）}\label{ux5047ux8bbe-2ux5364ux7d20ux63baux6742ux4f4dux70b9ux4e0e-n-ux578b-snse-ux8f7dux6d41ux5b50ux6d53ux5ea6ux53ca-zt-ux7684ux5173ux8054hypo_2}

\textbf{结论简述}：针对 n 型 SnSe 性能远落后 p 型 SnSe 的根因问题，提出以卤素（Cl/Br/I）取代 Se 位并结合 Sn 空位补偿来调控 n 型载流子浓度。假设掺杂剂的离子半径和缺陷形成能决定其固溶度及电离效率，进而控制电子浓度与晶格畸变程度。该假设可通过第一性原理缺陷化学计算预测各掺杂体系的缺陷形成能，再通过实验合成不同卤素掺杂浓度的 SnSe 多晶，结合变温 Hall 和输运测量检验。

\textbf{证据编号列表}： - \texttt{\cite{TE122}}（论文编号） - \texttt{\cite{TE115}}（论文编号） - \texttt{\cite{TE113}}（论文编号） - \texttt{{[}Novelty\ Verification{]}\ Overlap:\ none\ \textbar{}\ Novelty:\ 0.680\ (was\ 0.680)\ \textbar{}\ Queries:\ 3\ \textbar{}\ Results:\ 0}（新颖性查重）

\textbf{来源归属}： - 论文证据：\texttt{\cite{TE122}}、\texttt{\cite{TE115}}、\texttt{\cite{TE113}}（源自 hypotheses.json → evidence\_chain） - 新颖性查重：新颖性 0.680；查询 3 次；结果 0 条 - 数据库记录：无（hypotheses.json 中 external\_validation 为空）

\textbf{可追溯状态}：✅ 已溯源 ------ 证据编号在 discovery 产物（hypotheses.json → evidence\_chain、discovery\_report.md → Evidence Chain）中可逐层回查，且未被引用审计标记；论文编号对应的真实文献最终确认请回查检索缓存或人工复核。

\begin{center}\rule{0.5\linewidth}{0.5pt}\end{center}

\subsubsection{假设 3｜最优掺杂浓度随温度上移抑制双极效应的实验验证（hypo\_3）}\label{ux5047ux8bbe-3ux6700ux4f18ux63baux6742ux6d53ux5ea6ux968fux6e29ux5ea6ux4e0aux79fbux6291ux5236ux53ccux6781ux6548ux5e94ux7684ux5b9eux9a8cux9a8cux8bc1hypo_3}

\textbf{结论简述}：针对双极效应下最优掺杂随温度演化的实验验证缺失，提出在 Bi2Te3 基和 PbTe 材料中，通过制备一系列不同掺杂浓度的样品，系统测量 300--700 K 的电输运和热输运性质，绘制最优载流子浓度 n*(T) 相图。假设随着温度升高，本征激发增强，双极热导率和双极 Seebeck 抵消效应加剧，因此为了抑制双极效应，最优掺杂浓度应向更高值移动。 \textbf{数值验证结果} - \textbf{预测值（待实验验证）}: \texttt{10–30\%} --- 在 4 篇论文摘要中检索，未找到直接支撑。 \ldots{}

\textbf{证据编号列表}： - \texttt{\cite{TE082}}（论文编号） - \texttt{\cite{TE085}}（论文编号） - \texttt{\cite{TE090}}（论文编号） - \texttt{{[}Novelty\ Verification{]}\ Overlap:\ none\ \textbar{}\ Novelty:\ 0.740\ (was\ 0.740)\ \textbar{}\ Queries:\ 3\ \textbar{}\ Results:\ 2}（新颖性查重） - \texttt{预测值\ 10–30\%（待实验验证，建议方案：元素分析或TGA定量组成变化）} - \texttt{预测值\ 600\ K\ —\ 未在现有文献中找到直接支撑。建议通过系统性实验或第一性原理计算进行验证。} - \texttt{预测值\ 15–30\%（待实验验证，建议方案：元素分析或TGA定量组成变化）} - \texttt{{[}Overlap{]}\ "Thermoelectric\ performance\ of\ P-N-P\ abrupt\ heterostructures\ vertical\ to\ temperature\ gradient"\ (sim=0.048)} - \texttt{{[}Overlap{]}\ "Porosity-mediated\ High-performance\ Thermoelectric\ Materials"\ (sim=0.032)}

\textbf{来源归属}： - 论文证据：\texttt{\cite{TE082}}、\texttt{\cite{TE085}}、\texttt{\cite{TE090}}（源自 hypotheses.json → evidence\_chain） - 新颖性查重：新颖性 0.740；查询 3 次；结果 2 条 - 其他证据：预测值 10--30\%（待实验验证，建议方案：元素分析或TGA定量组成变化） - 其他证据：预测值 600 K --- 未在现有文献中找到直接支撑。建议通过系统性实验或第一性原理计算进行验证。 - 其他证据：预测值 15--30\%（待实验验证，建议方案：元素分析或TGA定量组成变化） - 其他证据：{[}Overlap{]} ``Thermoelectric performance of P-N-P abrupt heterostructures vertical to temperature gradient'' (sim=0.048) - 其他证据：{[}Overlap{]} ``Porosity-mediated High-performance Thermoelectric Materials'' (sim=0.032) - 数据库记录：无（hypotheses.json 中 external\_validation 为空）

\textbf{可追溯状态}：✅ 已溯源 ------ 证据编号在 discovery 产物（hypotheses.json → evidence\_chain、discovery\_report.md → Evidence Chain）中可逐层回查，且未被引用审计标记；论文编号对应的真实文献最终确认请回查检索缓存或人工复核。

\begin{center}\rule{0.5\linewidth}{0.5pt}\end{center}

\subsubsection{假设 4｜共振掺杂浓度与 DOS 局部异常及 Seebeck 增强的定量标度（hypo\_4）}\label{ux5047ux8bbe-4ux5171ux632fux63baux6742ux6d53ux5ea6ux4e0e-dos-ux5c40ux90e8ux5f02ux5e38ux53ca-seebeck-ux589eux5f3aux7684ux5b9aux91cfux6807ux5ea6hypo_4}

\textbf{结论简述}：针对共振掺杂缺乏定量标度规律的问题，提出在 half-Heusler HfZrCoSnSb 中，通过控制 Al 共振掺杂浓度（0--1.5 at\%）并利用第一性原理计算费米能级附近态密度（DOS）的局部畸变幅度，再与实验测量的 Seebeck 系数和功率因子关联。假设共振杂质产生的 DOS 局部异常幅度越大，Seebeck 系数的增强越显著，但超过临界浓度后杂质带扩展会破坏共振态，导致 Seebeck 骤降。

\textbf{证据编号列表}： - \texttt{\cite{TE146}}（论文编号） - \texttt{\cite{TE087}}（论文编号） - \texttt{\cite{TE101}}（论文编号） - \texttt{{[}Novelty\ Verification{]}\ Overlap:\ none\ \textbar{}\ Novelty:\ 0.820\ (was\ 0.820)\ \textbar{}\ Queries:\ 3\ \textbar{}\ Results:\ 0}（新颖性查重）

\textbf{来源归属}： - 论文证据：\texttt{\cite{TE146}}、\texttt{\cite{TE087}}、\texttt{\cite{TE101}}（源自 hypotheses.json → evidence\_chain） - 新颖性查重：新颖性 0.820；查询 3 次；结果 0 条 - 数据库记录：无（hypotheses.json 中 external\_validation 为空）

\textbf{可追溯状态}：✅ 已溯源 ------ 证据编号在 discovery 产物（hypotheses.json → evidence\_chain、discovery\_report.md → Evidence Chain）中可逐层回查，且未被引用审计标记；论文编号对应的真实文献最终确认请回查检索缓存或人工复核。

\begin{center}\rule{0.5\linewidth}{0.5pt}\end{center}

\subsubsection{假设 5｜Yb 填充方钴矿填充分数对晶格热导率与迁移率的定量影响及批次稳定性（hypo\_5）}\label{ux5047ux8bbe-5yb-ux586bux5145ux65b9ux94b4ux77ffux586bux5145ux5206ux6570ux5bf9ux6676ux683cux70edux5bfcux7387ux4e0eux8fc1ux79fbux7387ux7684ux5b9aux91cfux5f71ux54cdux53caux6279ux6b21ux7a33ux5b9aux6027hypo_5}

\textbf{结论简述}：针对填充方钴矿批次复现性差的问题，提出名义填充分数 y 与实际填充分数之间的偏差是导致 ZT 不一致的主要原因。假设 Yb 填充原子在方钴矿笼子中产生``rattling''效应，可定量降低晶格热导率；但过量填充会引入附加的电子散射并改变载流子浓度。通过制备一系列 y=0--0.4 的 Yb\_yFe4Sb12/Yb\_yCo4Sb12 样品，测定实际填充分数、晶格热导率、载流子迁移率和 ZT，可确定性能最优且对 y 不敏感的填充范围。 \textbf{数值验证结果} - **预测值（待实验\ldots{}

\textbf{证据编号列表}： - \texttt{\cite{TE048}}（论文编号） - \texttt{\cite{TE035}}（论文编号） - \texttt{\cite{TE037}}（论文编号） - \texttt{\cite{TE051}}（论文编号） - \texttt{{[}Novelty\ Verification{]}\ Overlap:\ none\ \textbar{}\ Novelty:\ 0.650\ (was\ 0.650)\ \textbar{}\ Queries:\ 3\ \textbar{}\ Results:\ 0}（新颖性查重） - \texttt{预测值\ 40–60\%（待实验验证，建议方案：元素分析或TGA定量组成变化）}

\textbf{来源归属}： - 论文证据：\texttt{\cite{TE048}}、\texttt{\cite{TE035}}、\texttt{\cite{TE037}}、\texttt{\cite{TE051}}（源自 hypotheses.json → evidence\_chain） - 新颖性查重：新颖性 0.650；查询 3 次；结果 0 条 - 其他证据：预测值 40--60\%（待实验验证，建议方案：元素分析或TGA定量组成变化） - 数据库记录：无（hypotheses.json 中 external\_validation 为空）

\textbf{可追溯状态}：✅ 已溯源 ------ 证据编号在 discovery 产物（hypotheses.json → evidence\_chain、discovery\_report.md → Evidence Chain）中可逐层回查，且未被引用审计标记；论文编号对应的真实文献最终确认请回查检索缓存或人工复核。

\section{热电材料 ZT 优化 --- 研究空白分析报告（Gap Report）}\label{gap:ux70edux7535ux6750ux6599-zt-ux4f18ux5316-ux7814ux7a76ux7a7aux767dux5206ux6790ux62a5ux544agap-report}

\textbf{主题}: Thermoelectric Materials ZT Optimization \textbf{分析基准}: 209 篇文献摘要（arxiv 160 + semantic\_scholar 49） \textbf{日期}: 2026-08-02

\begin{center}\rule{0.5\linewidth}{0.5pt}\end{center}

\subsection{Gap 1：ML 预测 ZT 与实验验证之间的系统性断层}\label{gap:gap-1ml-ux9884ux6d4b-zt-ux4e0eux5b9eux9a8cux9a8cux8bc1ux4e4bux95f4ux7684ux7cfbux7edfux6027ux65adux5c42}

\begin{itemize}
\tightlist
\item
  \textbf{类型}: 缺失连接
\item
  \textbf{严重程度}: 高
\item
  \textbf{置信度}: 0.92
\item
  \textbf{描述}: ML 模型在热电性质预测上报告 R²=0.90--0.98 的高精度（\cite{TE057}），理论预测 ZT 可达 3.0--7.2（\cite{TE062}/\cite{TE096}/\cite{TE009}/\cite{TE002}），但仅有极少数预测转化为实验验证的高 ZT 材料。核心障碍：小数据问题、交叉验证采样偏差（忽略化学家族的''隐藏层次结构''）、结构表示不充分、热力学相稳定性被忽视（\cite{TE057}）。
\item
  \textbf{证据}: \cite{TE057} (10.1016/j.aichem.2026.100113), \cite{TE058}, \cite{TE011}, \cite{TE062} (10.1063/1.3518078), \cite{TE096}, \cite{TE009}
\item
  \textbf{矛盾点}: 大量理论 ZT\textgreater3 的预测（Bi₂Te₃ 薄膜 7.2、PbSe/PbTe 单层 5.3、Graphdiyne 3.0）vs 实验纪录仅 \textasciitilde2.6（SnSe 单晶），预测-实验鸿沟 ≥2-3 倍
\item
  \textbf{验证方案}: 选取 ML 预测的高 ZT 候选（如超饱和 Si-Ge-P），进行实验合成+输运测量闭环验证；采用 PCA 基采样与主动学习循环替代随机 CV
\end{itemize}

\subsection{Gap 2：n 型 SnSe 性能远落后 p 型 SnSe 的根因未解}\label{gap:gap-2n-ux578b-snse-ux6027ux80fdux8fdcux843dux540e-p-ux578b-snse-ux7684ux6839ux56e0ux672aux89e3}

\begin{itemize}
\tightlist
\item
  \textbf{类型}: 缺失连接
\item
  \textbf{严重程度}: 高
\item
  \textbf{置信度}: 0.88
\item
  \textbf{描述}: p 型 SnSe 单晶 ZT=2.6 创纪录（\cite{TE113}），但 n 型 SnSe 多晶最高仅 \textasciitilde1.8（\cite{TE115}），且普遍 \textless1。根本障碍：本征缺陷化学（Sn 空位天然 p 型）、强各向异性输运、掺杂剂选择受限（\cite{TE122}）。
\item
  \textbf{证据}: \cite{TE122} (10.1016/j.isci.2025.114025), \cite{TE115}, \cite{TE113}, \cite{TE119}, \cite{TE117}
\item
  \textbf{验证方案}: 系统研究卤素（Cl/Br/I）与过渡金属掺杂的缺陷形成能-载流子浓度-ZT 关联；结合 DFT 缺陷化学计算与实验
\end{itemize}

\subsection{Gap 3：填充方钴矿 ZT 的批次复现性差}\label{gap:gap-3ux586bux5145ux65b9ux94b4ux77ff-zt-ux7684ux6279ux6b21ux590dux73b0ux6027ux5dee}

\begin{itemize}
\tightlist
\item
  \textbf{类型}: 矛盾结论
\item
  \textbf{严重程度}: 中
\item
  \textbf{置信度}: 0.75
\item
  \textbf{描述}: p 型 RᵧFe₄ 方钴矿同类材料不同实验室报告 ZT 0.67/0.70/0.89 显著不一致（\cite{TE048}）；Yb 溶解度温度依赖性复杂（\cite{TE035}），填充原子种类（Ce/Yb/In/S）对性能影响规律缺乏统一框架。
\item
  \textbf{证据}: \cite{TE048}, \cite{TE035} (10.1016/J.JMAT.2015.03.008), \cite{TE037}, \cite{TE051}
\item
  \textbf{验证方案}: 建立标准合成-测量协议；对同配方进行多实验室盲测；定量填充分数-晶格热导-迁移率的关系
\end{itemize}

\subsection{Gap 4：双极效应下最优掺杂随温度演化的实验验证缺失}\label{gap:gap-4ux53ccux6781ux6548ux5e94ux4e0bux6700ux4f18ux63baux6742ux968fux6e29ux5ea6ux6f14ux5316ux7684ux5b9eux9a8cux9a8cux8bc1ux7f3aux5931}

\begin{itemize}
\tightlist
\item
  \textbf{类型}: 缺失连接
\item
  \textbf{严重程度}: 中
\item
  \textbf{置信度}: 0.80
\item
  \textbf{描述}: 理论计算表明双极材料最优掺杂浓度在 600K 时比单极材料高 10--30\%（\cite{TE082}），且最优 Fermi 能级随工作温度移动；但实验上''温度-掺杂-功率因子''联合优化研究稀缺，多数材料仅做单温度点优化。
\item
  \textbf{证据}: \cite{TE082} (10.1007/s11664-018-06857-1), \cite{TE085}, \cite{TE090}
\item
  \textbf{验证方案}: 对 Bi₂Te₃ 与 PbTe 进行变温 Hall+输运联合测量，绘制 n*(T) 相图；与 \cite{TE136} 载流子散射相图结合
\end{itemize}

\subsection{Gap 5：half-Heusler 掺杂多重角色的散射相图普适性未验证}\label{gap:gap-5half-heusler-ux63baux6742ux591aux91cdux89d2ux8272ux7684ux6563ux5c04ux76f8ux56feux666eux9002ux6027ux672aux9a8cux8bc1}

\begin{itemize}
\tightlist
\item
  \textbf{类型}: 未探索空间
\item
  \textbf{严重程度}: 中
\item
  \textbf{置信度}: 0.72
\item
  \textbf{描述}: ZrNiSn 基 half-Heusler 建立了载流子散射相图（化学掺杂同时屏蔽电离杂质/晶界/极性光学声子散射，但对 κₗ 影响可忽略，\cite{TE136}）；该相图框架是否适用于其他 half-Heusler 家族（TiNiSn、FeNbSb）或其他材料体系（方钴矿、PbTe）未验证。
\item
  \textbf{证据}: \cite{TE136} (10.1038/s41467-020-16913-2), \cite{TE146}, \cite{TE200}
\item
  \textbf{验证方案}: 在 FeNb₀.₈Ti₀.₂Sb 和 Yb 填充方钴矿中复现''掺杂类型-散射机制-κₗ不变性''的关联
\end{itemize}

\subsection{Gap 6：共振掺杂（resonant doping）的 DOS 局部异常-塞贝克增强定量标度缺失}\label{gap:gap-6ux5171ux632fux63baux6742resonant-dopingux7684-dos-ux5c40ux90e8ux5f02ux5e38-ux585eux8d1dux514bux589eux5f3aux5b9aux91cfux6807ux5ea6ux7f3aux5931}

\begin{itemize}
\tightlist
\item
  \textbf{类型}: 未探索空间
\item
  \textbf{严重程度}: 中
\item
  \textbf{置信度}: 0.70
\item
  \textbf{描述}: \textless1at\% Al 共振掺杂使 HfZrCoSnSb 的 PF 提升 65\%、ZT 0.8→1.5（\cite{TE146}），机制归于 Fermi 能级附近 DOS 局部异常；但''共振杂质浓度-DOS 异常幅度-Seebeck 增益''之间缺乏定量标度律，无法指导其他体系推广。
\item
  \textbf{证据}: \cite{TE146}, \cite{TE087} (PbTe K/Na 共掺), \cite{TE101}
\item
  \textbf{验证方案}: 对不同共振杂质（Al/Ga/In/Tl）在不同母体中的 DOS 异常幅度进行第一性原理定量计算+实验对照
\end{itemize}

\begin{center}\rule{0.5\linewidth}{0.5pt}\end{center}

\subsection{Gap 7：共振掺杂文献缺乏 Seebeck 定量数值（量化建模数据断层）--- 2026-08-11 新增}\label{gap:gap-7ux5171ux632fux63baux6742ux6587ux732eux7f3aux4e4f-seebeck-ux5b9aux91cfux6570ux503cux91cfux5316ux5efaux6a21ux6570ux636eux65adux5c42-2026-08-11-ux65b0ux589e}

\begin{itemize}
\tightlist
\item
  \textbf{类型}: 缺失连接
\item
  \textbf{严重程度}: 中
\item
  \textbf{置信度}: 0.65
\item
  \textbf{描述}: 量化验证阶段发现：half-Heusler 共振掺杂文献（\cite{TE146}: HfZrCoSnSb+Al PF↑65\%、\cite{TE087}: PbTe K/Na 共掺、\cite{TE101}: Cr 共振掺杂）报告了 ZT 提升与功率因子增益，但摘要层面几乎不报告 Seebeck 系数（µV/K）的定量数值。对假设 3（共振掺杂浓度-DOS-Seebeck 定量标度）执行 run\_model\_comparison 与 symbolic\_regression 时，因文献中无 Seebeck 数值点（µV/K 单位桶匹配 0 个数据）而无法建模------``共振掺杂浓度-Seebeck 增益''的定量标度（Gap 6 的核心诉求）受制于文献数据报告格式，无法从摘要级证据链直接验证。
\item
  \textbf{证据}: \cite{TE146}, \cite{TE087}, \cite{TE101}（均无 Seebeck µV/K 数值）；量化验证结果：假设 3 数据点=0
\item
  \textbf{验证方案}: ① 从全文（非摘要）提取 half-Heusler 变温 Seebeck 数据（如 HfZrCoSnSb 的 S(T) 曲线）；② 对 HfZrCoSnSb:Al 系列（0-1.5 at\%）进行第一性原理输运计算（Boltzmann/BTE），补全''浓度-Seebeck''数据点；③ 推动领域数据标准：在摘要中报告峰值 Seebeck 与功率因子数值
\end{itemize}

\subsection{Gap 1 更新（2026-08-11 量化验证强化）：}\label{gap:gap-1-ux66f4ux65b02026-08-11-ux91cfux5316ux9a8cux8bc1ux5f3aux5316}

\begin{itemize}
\tightlist
\item
  \textbf{置信度}: 0.92 → 0.93
\item
  \textbf{新增证据}: 跨材料温度-ZT 单变量建模（15-17 点）候选三次 R²=0.26-0.29、嵌套 F 检验不显著（p=0.14-0.16）、符号回归仅得过拟合复杂表达式 → ZT 无法由单一温度变量预测，必须引入材料描述符（带隙、有效质量、声速、德拜温度）；且外部数据库（MP/OQMD/NOMAD）无 ZT 实验数据 → 数据层断层确认
\item
  \textbf{验证方案补充}: 构建多描述符（T + 带隙 + 有效质量 + 声速）ZT 回归模型，检验 R² 是否显著提升
\end{itemize}

\subsection{Gap 优先级排序（用于假设生成）}\label{gap:gap-ux4f18ux5148ux7ea7ux6392ux5e8fux7528ux4e8eux5047ux8bbeux751fux6210}

\begin{enumerate}
\def\labelenumi{\arabic{enumi}.}
\tightlist
\item
  \textbf{Gap 1}（ML-实验断层）--- 严重程度：高 --- 置信度 0.92
\item
  \textbf{Gap 2}（n 型 SnSe 落后）--- 严重程度：高 --- 置信度 0.88
\item
  \textbf{Gap 6}（共振掺杂定量标度）--- 严重程度：中 --- 置信度 0.70
\item
  \textbf{Gap 4}（双极掺杂温度演化）--- 严重程度：中 --- 置信度 0.80
\item
  \textbf{Gap 3}（方钴矿复现性）--- 严重程度：中 --- 置信度 0.75
\item
  \textbf{Gap 5}（散射相图普适性）--- 严重程度：中 --- 置信度 0.72
\end{enumerate}

\printbibliography[title=参考文献]
\end{document}
